% GENERATED FILE. Edit canonical lessons, then run npm run generate:books.
% canonical-source-hash: fnv1a64:b46ca77f2ce523ab
% canonical-lessons: JA-C460-mokuhyou, JA-C460-seikou-suru, JA-C460-jitsuryoku, JA-C460-chousen-suru, JA-C460-tenshoku-suru

\chapter{Goals and Success}
\label{ch:ja-goals-and-success}
\hlchaptermodality{460}

\begin{chapteropening}
\textbf{By the end of this chapter:} \emph{I can talk about goals, success and challenges.}

Complete the last lesson of chapter 460: I can talk about goals, success and challenges.
\end{chapteropening}

\section[\texorpdfstring{\ja{もくひょう} (mokuhyō)}{もくひょう (mokuhyō)}]{\ja{もくひょう} (mokuhyō) — a goal; a target}

\label{lesson:JA-C460-mokuhyou}



\begin{quote}
Before the new one: say the Japanese for one's field of study, then the Japanese for a career path.
\end{quote}

\subsection*{You'll want to know: \ja{もくひょう}}
\textbf{\ja{もくひょう}} — \emph{mokuhyō} — \textquotedblleft{}a goal; a target\textquotedblright{}.

\begin{quote}
\textbf{\ja{ことしの} \ja{もくひょうは} \ja{しあいで} \ja{かつ} \ja{ことです。}} — \emph{kotoshi no mokuhyō wa shiai de katsu koto desu} — \textquotedblleft{}My goal this year is to win a match.\textquotedblright{}
\end{quote}

\begin{cousinweb}
\textbf{\ja{もくひょう}} (\emph{mokuhyō}) is Sino-Japanese (kango): two Chinese characters read with their on'yomi, readings borrowed from Chinese. They are \emph{moku}, \textquotedblleft{}eye\textquotedblright{} (Middle Chinese \emph{mjuwk}) and \emph{hyō}, \textquotedblleft{}mark\textquotedblright{} (\emph{pjiew}) (KANJIDIC2; Middle Chinese in Baxter's transcription).
\end{cousinweb}

\begin{grammarlens}[title={What you've built}]
One more word for describing what you lived and felt.
\end{grammarlens}

\subsection*{Your turn}
\begin{itemize}
\raggedright
  \item \emph{Say it:} \emph{mokuhyō}
  \item \emph{Say it:} \emph{mokuhyō}, once more
  \item \emph{Say it:} \emph{kotoshi no mokuhyō wa shiai de katsu koto desu} — My goal this year is to win a match.
  \item \emph{Recall:} say the Japanese for one's field of study, then the Japanese for a career path, then say \emph{mokuhyō} again
\end{itemize}

\begin{tcolorbox}[breakable,colback=teal!4,colframe=teal!35!black,title={Before you move on}]
What does \ja{もくひょう} mean? (\textquotedblleft{}A goal; a target\textquotedblright{}.) Say it once more.
\end{tcolorbox}
\section[\texorpdfstring{\ja{せいこうする} (seikō suru)}{せいこうする (seikō suru)}]{\ja{せいこうする} (seikō suru) — to succeed}

\label{lesson:JA-C460-seikou-suru}



\begin{quote}
Before the new one: say the Japanese for a career path, then the Japanese for a goal.
\end{quote}

\subsection*{You'll want to know: \ja{せいこうする}}
\textbf{\ja{せいこうする}} — \emph{seikō suru} — \textquotedblleft{}to succeed\textquotedblright{}.

\begin{quote}
\textbf{\ja{まいにち} \ja{れんしゅうして、せいこうしました。}} — \emph{mainichi renshū shite, seikō shimashita} — \textquotedblleft{}I practised every day and succeeded.\textquotedblright{}
\end{quote}

Polite \textbf{\ja{せいこうします}}, \emph{seikōshimasu}; a \ja{する} verb. Past (polite): \textbf{\ja{あにが} \ja{せいこうしました。}}, \emph{ani ga seikō shimashita}.

\begin{cousinweb}
\textbf{\ja{せいこうする}} (\emph{seikō suru}) is Sino-Japanese (kango) \emph{seikō} with the native verb \emph{suru}, \textquotedblleft{}to do\textquotedblright{}: two Chinese characters read with their on'yomi, readings borrowed from Chinese. They are \emph{sei}, \textquotedblleft{}become\textquotedblright{} (Middle Chinese \emph{dzyeng}) and \emph{kō}, \textquotedblleft{}success\textquotedblright{} (\emph{kuwng}) (KANJIDIC2; Middle Chinese in Baxter's transcription).
\end{cousinweb}

\begin{grammarlens}[title={What you've built}]
One more word for describing what you lived and felt.
\end{grammarlens}

\subsection*{Your turn}
\begin{itemize}
\raggedright
  \item \emph{Say it:} \emph{seikō suru}
  \item \emph{Say it:} \emph{seikō suru}, once more
  \item \emph{Say it:} \emph{mainichi renshū shite, seikō shimashita} — I practised every day and succeeded.
  \item \emph{Recall:} say the Japanese for a career path, then the Japanese for a goal, then say \emph{seikō suru} again
\end{itemize}

\begin{tcolorbox}[breakable,colback=teal!4,colframe=teal!35!black,title={Before you move on}]
What does \ja{せいこうする} mean? (\textquotedblleft{}To succeed\textquotedblright{}.) Say it once more.
\end{tcolorbox}
\section[\texorpdfstring{\ja{じつりょく} (jitsuryoku)}{じつりょく (jitsuryoku)}]{\ja{じつりょく} (jitsuryoku) — real ability}

\label{lesson:JA-C460-jitsuryoku}



\begin{quote}
Before the new one: say the Japanese for a goal, then the Japanese for to succeed.
\end{quote}

\subsection*{You'll want to know: \ja{じつりょく}}
\textbf{\ja{じつりょく}} — \emph{jitsuryoku} — \textquotedblleft{}real ability\textquotedblright{}.

\begin{quote}
\textbf{\ja{あの} \ja{せんしゅは} \ja{じつりょくが} \ja{あります。}} — \emph{ano senshu wa jitsuryoku ga arimasu} — \textquotedblleft{}That player has real ability.\textquotedblright{}
\end{quote}

\begin{cousinweb}
\textbf{\ja{じつりょく}} (\emph{jitsuryoku}) is Sino-Japanese (kango): two Chinese characters read with their on'yomi, readings borrowed from Chinese. They are \emph{jitsu}, \textquotedblleft{}reality\textquotedblright{} (Middle Chinese \emph{zyit}) and \emph{ryoku}, \textquotedblleft{}power\textquotedblright{} (\emph{lik}) (KANJIDIC2; Middle Chinese in Baxter's transcription).
\end{cousinweb}

\begin{grammarlens}[title={What you've built}]
One more word for describing what you lived and felt.
\end{grammarlens}

\subsection*{Your turn}
\begin{itemize}
\raggedright
  \item \emph{Say it:} \emph{jitsuryoku}
  \item \emph{Say it:} \emph{jitsuryoku}, once more
  \item \emph{Say it:} \emph{ano senshu wa jitsuryoku ga arimasu} — That player has real ability.
  \item \emph{Recall:} say the Japanese for a goal, then the Japanese for to succeed, then say \emph{jitsuryoku} again
\end{itemize}

\begin{tcolorbox}[breakable,colback=teal!4,colframe=teal!35!black,title={Before you move on}]
What does \ja{じつりょく} mean? (\textquotedblleft{}Real ability\textquotedblright{}.) Say it once more.
\end{tcolorbox}
\section[\texorpdfstring{\ja{ちょうせんする} (chōsen suru)}{ちょうせんする (chōsen suru)}]{\ja{ちょうせんする} (chōsen suru) — to take on a challenge}

\label{lesson:JA-C460-chousen-suru}



\begin{quote}
Before the new one: say the Japanese for to succeed, then the Japanese for real ability.
\end{quote}

\subsection*{You'll want to know: \ja{ちょうせんする}}
\textbf{\ja{ちょうせんする}} — \emph{chōsen suru} — \textquotedblleft{}to take on a challenge\textquotedblright{}.

\begin{quote}
\textbf{\ja{むすこは} \ja{たかい} \ja{やまに} \ja{ちょうせんしたいです。}} — \emph{musuko wa takai yama ni chōsen shitai desu} — \textquotedblleft{}My son wants to take on a high mountain.\textquotedblright{}
\end{quote}

Polite \textbf{\ja{ちょうせんします}}, \emph{chōsenshimasu}; a \ja{する} verb. Past (polite): \textbf{\ja{せんしゅが} \ja{ちょうせんしました。}}, \emph{senshu ga chōsen shimashita}.

\begin{cousinweb}
\textbf{\ja{ちょうせんする}} (\emph{chōsen suru}) is Sino-Japanese (kango) \emph{chōsen} with the native verb \emph{suru}, \textquotedblleft{}to do\textquotedblright{}: two Chinese characters read with their on'yomi, readings borrowed from Chinese. They are \emph{chō}, \textquotedblleft{}challenge\textquotedblright{} (Middle Chinese \emph{thew}) and \emph{sen}, \textquotedblleft{}battle\textquotedblright{} (\emph{tsyenH}) (KANJIDIC2; Middle Chinese in Baxter's transcription).
\end{cousinweb}

\begin{grammarlens}[title={What you've built}]
One more word for describing what you lived and felt.
\end{grammarlens}

\subsection*{Your turn}
\begin{itemize}
\raggedright
  \item \emph{Say it:} \emph{chōsen suru}
  \item \emph{Say it:} \emph{chōsen suru}, once more
  \item \emph{Say it:} \emph{musuko wa takai yama ni chōsen shitai desu} — My son wants to take on a high mountain.
  \item \emph{Recall:} say the Japanese for to succeed, then the Japanese for real ability, then say \emph{chōsen suru} again
\end{itemize}

\begin{tcolorbox}[breakable,colback=teal!4,colframe=teal!35!black,title={Before you move on}]
What does \ja{ちょうせんする} mean? (\textquotedblleft{}To take on a challenge\textquotedblright{}.) Say it once more.
\end{tcolorbox}
\section[\texorpdfstring{\ja{てんしょくする} (tenshoku suru)}{てんしょくする (tenshoku suru)}]{\ja{てんしょくする} (tenshoku suru) — to change jobs}

\label{lesson:JA-C460-tenshoku-suru}



\begin{quote}
Before the new one: say the Japanese for real ability, then the Japanese for to take on a challenge.
\end{quote}

\subsection*{You'll want to know: \ja{てんしょくする}}
\textbf{\ja{てんしょくする}} — \emph{tenshoku suru} — \textquotedblleft{}to change jobs\textquotedblright{}.

\begin{quote}
\textbf{\ja{あには} \ja{きょねん} \ja{てんしょくしました。}} — \emph{ani wa kyonen tenshoku shimashita} — \textquotedblleft{}My elder brother changed jobs last year.\textquotedblright{}
\end{quote}

Polite \textbf{\ja{てんしょくします}}, \emph{tenshokushimasu}; a \ja{する} verb. Past (polite): \textbf{\ja{ちちが} \ja{てんしょくしました。}}, \emph{chichi ga tenshoku shimashita}.

\begin{cousinweb}
\textbf{\ja{てんしょくする}} (\emph{tenshoku suru}) is Sino-Japanese (kango) \emph{tenshoku} with the native verb \emph{suru}, \textquotedblleft{}to do\textquotedblright{}: two Chinese characters read with their on'yomi, readings borrowed from Chinese. They are \emph{ten}, \textquotedblleft{}change\textquotedblright{} (Middle Chinese \emph{trjwenX}) and \emph{shoku}, \textquotedblleft{}employment\textquotedblright{} (\emph{tsyik}) (KANJIDIC2; Middle Chinese in Baxter's transcription).
\end{cousinweb}

\begin{grammarlens}[title={What you've built}]
One more word for describing what you lived and felt.
\end{grammarlens}

\subsection*{Your turn}
\begin{itemize}
\raggedright
  \item \emph{Say it:} \emph{tenshoku suru}
  \item \emph{Say it:} \emph{tenshoku suru}, once more
  \item \emph{Say it:} \emph{ani wa kyonen tenshoku shimashita} — My elder brother changed jobs last year.
  \item \emph{Recall:} say the Japanese for real ability, then the Japanese for to take on a challenge, then say \emph{tenshoku suru} again
\end{itemize}

\begin{tcolorbox}[breakable,colback=teal!4,colframe=teal!35!black,title={Before you move on}]
What does \ja{てんしょくする} mean? (\textquotedblleft{}To change jobs\textquotedblright{}.) Say it once more.
\end{tcolorbox}
